\section*{الفصل \ref{ch:b2:population-genetics} --- وراثة الجمهرات}

\begin{solution}{exo:b2:population-genetics:1}
$p(M) = (1280 + 320)/2000 = 0.8$، و$q(N) = 0.2$. والمنتظر من $MM$:
$0.64\times 1000 = 640$؛ ومن $MN$: $2\times 0.8\times 0.2\times 1000 =
320$؛ ومن $NN$: 40 --- وهي الأعداد الملحوظة بالضبط. فالجمهرة
في توازن عند هذا الموقع.
\end{solution}

\begin{solution}{exo:b2:population-genetics:2}
$q = \sqrt{1/20000} = 0.0071$؛ والحاملون $2pq \approx 0.014$، أي شخص
من كل 70. والحاملون يفوقون المصابين بنسبة $2p/q \approx
280$ إلى واحد.
\end{solution}

\begin{solution}{exo:b2:population-genetics:3}
\omterm{def:b2:population-genetics:fitness}{اللياقة}: عدد خلف \omterm{def:b2:meiosis-heredity:vocabulary}{نمط وراثي} المنتظر نسبةً إلى
الأفضل. \omterm{def:b2:population-genetics:fitness}{ومعامل الانتقاء} $s$: عجز \omterm{def:b2:population-genetics:fitness}{اللياقة}، $w = 1 -
s$. \omterm{thm:b2:population-genetics:heritability}{وقابلية التوريث} $h^{2}$: كسر التباين الظاهري
الذي هو وراثي جمعي، أي ميل انحدار الخلف على
الآباء. والقوى: \omterm{def:b2:genome-diversification:mutation}{الطفرة}، والانتقاء، والانجراف، والهجرة، والتزاوج
غير العشوائي (والأخير يغير تواترات الأنماط الوراثية لا تواترات الأليلات).
\end{solution}

\begin{solution}{exo:b2:population-genetics:4}
الموجّه: تسوّد الفراشة الفلفلية، وعامله الطيور التي تصطاد
بالنظر على جذوع سوّدها السخام. والمثبِّت: وزن ولادة الإنسان،
وعامله وفيات المواليد الصغار جدًا والكبار جدًا. والموازن:
الخلية المنجلية، وعاملاه الملاريا (ضد $AA$) وفقر الدم (ضد
$SS$) معًا.
\end{solution}

\begin{solution}{exo:b2:population-genetics:5}
$q_t = q_0/(1 + tq_0)$: فالتنصيف يستغرق $1/q_0 = 50$ جيلًا؛
وبلوغ $0.001$ يستغرق $1/0.001 - 1/0.02 = 950$ جيلًا. فجل
نسخ \omterm{def:b2:meiosis-heredity:vocabulary}{الأليل} تجلس في متغايري اللواقح، خفيةً على
الانتقاء؛ ومنع المصابين من التكاثر (وهو ما تفعله الطبيعة
أصلًا، إذ هم موتى) لا يغير شيئًا يُقاس في
عمر إنسان.
\end{solution}

\begin{solution}{exo:b2:population-genetics:6}
عند $p = 0.3$: $\bar w = 0.09 + 0.42 + 0.8\times 0.49 = 0.902$؛ و$p' =
(0.09 + 0.21)/0.902 = 0.333$؛ و$\Delta p = +0.033$. وعند $p = 0.9$: $\bar w
= 0.998$، و$p' = 0.9018$، و$\Delta p = +0.0018$. فعند $p = 0.9$ ليس من الجمهرة
إلا $q^{2} = 0.01$ من نمط $aa$: \omterm{def:b2:meiosis-heredity:vocabulary}{فالأليل} الذي يفعل الانتقاء
ضده مختبئ في متغايري اللواقح، و$\Delta p = spq^{2}/\bar w$
يحمل العامل $q^{2}$.
\end{solution}

\begin{solution}{exo:b2:population-genetics:7}
عند $s = 0.12$ ضد $AA$ و$t = 0.8$ ضد $SS$ يكون $\hat q =
s/(s + t) = 0.12/0.92 = 0.13$. والمواليد المصابون $\hat q^{2} = 0.017$،
أي نحو طفل من كل 60 --- وهو ثمن الحماية التي
ينعم بها متغايرو اللواقح.
\end{solution}

\begin{solution}{exo:b2:population-genetics:8}
المتنحي المميت: $\hat q = \sqrt{\mu} = \sqrt{2\times 10^{-5}} =
0.0045$؛ والمواليد المصابون $\hat q^{2} = \mu = 2\times 10^{-5}$، أي واحد من
\num{50000}. والمميت \omterm{def:b2:meiosis-heredity:vocabulary}{السائد}: تُزال كل نسخة في
الجيل الذي تظهر فيه، فيكون المواليد المصابون هم الطفرات
الجديدة وحدها، $2\mu = 4\times 10^{-5}$ (بعرسين لكل مولود)، أي واحد من
\num{25000} --- فالمتنحي يخبئ أليلاته بنسبة مئتين إلى واحد،
\omterm{def:b2:meiosis-heredity:vocabulary}{والسائد} لا يخبئ شيئًا.
\end{solution}

\begin{solution}{exo:b2:population-genetics:9}
$H_t = 0.5\,(1 - 1/100)^{t}$: فبعد 10 أجيال $0.45$؛ وبعد 50،
$0.30$؛ وبعد 200، $0.067$. ولأجل $(1 - 1/2N)^{100} = 0.9$:
يكون $100/(2N) \approx -\ln 0.9 = 0.105$، فيكون $N \approx 475$، أي نحو 500
فرد متكاثر.
\end{solution}

\begin{solution}{exo:b2:population-genetics:10}
عشرون نسخة مورّثية: $P(\text{الغياب}) = 0.9^{20} = 0.12$. والنسخة
الواحدة بين 20 تواترها $0.05$، \omterm{def:b2:meiosis-heredity:vocabulary}{والأليل} المحايد يُثبَّت
باحتمال يساوي تواتره: أي $0.05$. وفي بر رئيس فيه، مثلًا،
\num{100000} طائر يكون لنسخة واحدة احتمال $5\times 10^{-6}$
للتثبيت. فالجزيرة مكان تضيع فيه الأليلات النادرة و
يستطيع الباقي منها الاستيلاء: فالانجراف يفقر ويفرّق معًا.
\end{solution}

\begin{solution}{exo:b2:population-genetics:11}
$R = h^{2}S = 0.3\times 1.4\times \qty{1000}{L} = \qty{420}{L}$ في
الجيل. وبثيران من أعلى \qty{1}{\%} يكون الفرق
المتوسط $(1.4 + 2.7)/2 = 2.05$ انحراف معياري ويكون $R =
0.3\times 2050 = \qty{615}{L}$. وتبطؤ الاستجابة لأن الانتقاء
يستهلك التباين الجمعي --- إذ تمضي الأليلات المحبَّذة إلى التثبيت
وتهبط $h^{2}$ --- ولأن التغيرات المرتبطة في صفات
أخرى (الخصوبة والصحة) تفرض كلفًا لا يراها الانتقاء على
المردود وحده.
\end{solution}

\begin{solution}{exo:b2:population-genetics:12}
الانتقاء يرى الأنماط الظاهرية. \omterm{def:b2:meiosis-heredity:vocabulary}{والأليل} المتنحي في \omterm{def:b2:meiosis-heredity:vocabulary}{متغاير اللواقح}
لا يصنع نمطًا ظاهريًا، فلا يستطيع الانتقاء إزالته، ويهبط
تواتره على صورة $1/t$، أبطأ فأبطأ. \omterm{def:b2:genome-diversification:mutation}{والطفرة} المحايدة
لا تصنع فرقًا في \omterm{def:b2:population-genetics:fitness}{اللياقة}؛ ومصيرها للانجراف وحده، وهي
تُثبَّت باحتمال $1/2N$. والصفة التي قابلية توريثها صفر
قد تخضع لانتقاء شديد --- فقد يختلف الناجون كثيرًا عن
الجمهرة --- ومع ذلك لا يتغير الجيل التالي، لأن
الفروق لم تُورَّث. فالانتقاء لا يستطيع أن يفعل إلا في
فروق موروثة تصنع فرقًا في \omterm{def:b2:population-genetics:fitness}{اللياقة}؛ وكل ما عداها
يتطور بالصدفة أو لا يتطور.
\end{solution}

\begin{solution}{pb:b2:population-genetics:1}
\textbf{1.} كسر السود $1 - q^{2} = 1 - 0.999^{2} = 0.002$،
أي فراشة من كل 500؛ وكسر أليلات $M$ في متغايري اللواقح هو
$2pq/(2pq + 2p^{2}) = q = 0.999$: أي كلها عمليًا.
\textbf{2.} عند $\bar w = 1 - sq^{2}$، تكون نسخ \omterm{def:b2:meiosis-heredity:vocabulary}{الأليل} الفاتح
بعد الانتقاء $pq\cdot 1 + q^{2}(1 - s)$ مقسومة على $\bar w$، فيكون $q'
= (pq + q^{2} - sq^{2})/(1 - sq^{2}) = q(1 - sq)/(1 - sq^{2})$.
\textbf{3.} $q_1 = 0.999\times(1 - 0.2997)/(1 - 0.2994) = 0.99857$:
أي تغير قدره $-0.0014$.
\textbf{4.} $q = \sqrt{0.1} = 0.32$.
\textbf{5.} العلاقة التراجعية، إذا كُرّرت، تستغرق 28 جيلًا عند $s =
0.3$؛ والتقريب يعطي الرتبة نفسها (فالتغير
$-0.0003$ أولًا، ثم يتسارع إذ يكبر $p$). أما الأجيال 50
الملحوظة فتوافق $s$ أصغر، نحو $0.2$ (إذ تعطي العلاقة
التراجعية 45 جيلًا عند $s = 0.2$): أي انتقاء بمقدار \qtyrange{20}{30}{\%}
في الجيل، وهو هائل بمقاييس معظم الأليلات.
\textbf{6.} عند $w_{MM} = w_{Mm} = 1 - s$ و$w_{mm} = 1$: يكون $\bar w = 1
- s(1 - q^{2})$ و$p' = p(1 - s)/\bar w$، ومن ثم $\Delta p = p[(1 - s) -
\bar w]/\bar w = -spq^{2}/\bar w$. والشكل الفاتح لا يُحبَّذ إلا بوصفه
\omterm{def:b2:meiosis-heredity:vocabulary}{متماثل اللواقح} $mm$، وهو كسر $q^{2} = 0.1$ في البداية: فتبدأ
العودة بطيئة. وإذ يؤول $q \to 1$ يصير $\Delta p \approx -sp$ ويهبط $M$
هبوطًا أسّيًا، بعامل $1 - s$ في كل جيل: \omterm{def:b2:meiosis-heredity:vocabulary}{فالأليل}
\omterm{def:b2:meiosis-heredity:vocabulary}{السائد} لا مخبأ له. وكان الصعود صورةً في مرآة:
فيكون $\Delta p \approx +sp$ ما دام $M$ نادرًا (بداية أسّية)، ثم
$\Delta q \approx -sq^{2}p \to$ زحف إذ اختبأ \omterm{def:b2:meiosis-heredity:vocabulary}{الأليل} الفاتح في
متغايري اللواقح. والعلاقة التراجعية تعطي 27 جيلًا للعودة
من $p = 0.68$ إلى $0.001$، منها 16 لإنزال \omterm{def:b2:meiosis-heredity:vocabulary}{الأليل} إلى \qty{5}{\%}.
\textbf{7.} $24$ نسخة مورّثية: $0.95^{24} = 0.29$.
\textbf{8.} $H_1 = 0.4\,(1 - 1/24) = 0.383$.
\textbf{9.} $H_{100} = 0.383\,(1 - 1/80)^{100} = 0.383\times 0.284 =
0.109$: أي \qty{27}{\%} من قيمة البر الرئيس $0.4$.
\textbf{10.} $P(\text{التثبيت}) = 1/(2N) = 1/80 = 0.0125$؛ والزمن
المنتظر للتثبيت $4N = 160$ جيلًا.
\textbf{11.} دون \omterm{prop:b2:population-genetics:inbreeding}{تزاوج داخلي} يكون $q^{2} = 0.0025$؛ ومع $F = 0.3$ يكون
$q^{2} + Fpq = 0.0025 + 0.3\times 0.95\times 0.05 = 0.017$: أي سبعة
أمثال من الطيور المصابة.
\textbf{12.} مهاجر واحد في الجيل ($m = 1/40$) يكفي
لمنع الجزيرة من التباعد عند المواقع المحايدة و
لتجديد الأليلات التي أضاعها الانجراف؛ وهو يستورد أيضًا
أليلات البر الرئيس الضارة بتواترها هناك، وهي،
تحت تزاوج الجزيرة الداخلي، تُكشف متماثلةَ اللواقح --- فالحمل
يقرره ميزان الهجرة الوافدة والكشف.
\textbf{13.} $\hat q = \sqrt{10^{-5}} = 0.0032$؛ والمواليد المصابون $\hat
q^{2} = 10^{-5}$.
\textbf{14.} $\hat q = \sqrt{10^{-5}/0.1} = 0.01$؛ والمواليد المصابون
$10^{-4}$: فانتقاء أضعف بعشرة أضعاف يجعل \omterm{def:b2:meiosis-heredity:vocabulary}{الأليل} أشيع بثلاثة
أمثال والمرض أشيع بعشرة.
\textbf{15.} $\hat p = \mu/(hs) = 10^{-5}/0.05 = 2\times 10^{-4}$؛
والحاملون $2\hat p = 4\times 10^{-4}$.
\textbf{16.} من $\hat q = 0.0032$ إلى $0.01$: يرتفع \omterm{def:b2:meiosis-heredity:vocabulary}{الأليل} بمعدل
تقرره \omterm{def:b2:genome-diversification:mutation}{الطفرة}، $\mu = 10^{-5}$ في الجيل، فيستغرق
الاقتراب رتبة $\Delta q/\mu \approx 700$ جيل ---
أي \num{20000} سنة: فأي قرار يُتخذ الآن تحسه جمهرة
لن تشبه جمهرتنا في شيء.
\textbf{17.} يسهم كل موقع بعدد $2\hat q = 0.0063$ نسخة مميتة
في الشخص؛ وعلى \num{20000} موقع يكون ذلك نحو $130$ --- لكن ذلك يفترض
أن كل مورّثة تستطيع أن تطفر إلى \omterm{def:b2:meiosis-heredity:vocabulary}{متنحٍّ} مميت؛ وبالتقدير
المألوف البالغ بضعة آلاف مورّثة أساسية، يحمل الشخص
من عدة إلى بضع عشرات من المكافئات المميتة في حالة تغاير اللواقح.
\textbf{18.} كل منها عند موقع مختلف وكل منها نادر، فاحتمال
أن يحمل زوجان عشوائيان الأليلَ نفسه صغير؛
وأبناء العم يشتركون في ثُمن مورّثاتهم بالنسب، فلكل مميت من
مميتات ابن العم العديدة يكون خطر أن يكون الطفل \omterm{def:b2:meiosis-heredity:vocabulary}{متماثل اللواقح}
$1/16$ لكل مميت مشترك --- ووفيات الأطفال الزائدة في زواج
أبناء العم مقيسة ببضعة في المئة.
\textbf{19.} $S = 1.16\times \qty{800}{L} = \qty{930}{L}$؛ و$R = 0.4\times
930 = \qty{370}{L}$.
\textbf{20.} خمسة أجيال: نحو \qty{1860}{L}؛ وهو أصغر في
الواقع لأن التباين الجمعي يُستنفد، ولأن $h^{2}$
قُدّرت في القطيع الأصلي، ولأن البيئة
(العلف والإيواء) التي صنعت البقرات الأفضل قد لا تُنقل.
\textbf{21.} الفرق المتوسط $(1.16 + 2.4)/2 = 1.78$ انحراف
معياري، و$S = \qty{1424}{L}$، و$R = 0.4\times 1424 = \qty{570}{L}$.
\textbf{22.} المستنسخات بلا تباين وراثي، فيكون $h^{2} = 0$ ويكون $R =
0$ مهما كان $S$: \omterm{thm:b2:population-genetics:heritability}{فقابلية التوريث} خاصة لجمهرة، أي
حصة تباينها الوراثي، لا خاصة للصفة.
\textbf{23.} $R = 0.7\times 0.6 = \qty{0.42}{mm}$ أعمق.
\textbf{24.} الاستجابة نفسها $\qty{370}{L}$ تحتاج $S = 370/0.1 =
\qty{3700}{L}$، أي $4.6$ انحراف معياري --- وذلك بإبقاء أقل
من بقرة واحدة من كل عشرة آلاف: وهو غير عملي، ولذلك
تُحسَّن الصفات ضعيفة التوريث ببطء أو بوسائل أخرى
(اختبار النسل، والعائلات الأكبر).
\textbf{25.} الفراشة: $s \approx 0.2$ إلى $0.3$، و28 جيلًا عند $0.3$
و45 عند $0.2$، في مقابل 50 ملحوظة؛ والجزيرة: $H_{100} \approx 0.11$،
أي \qty{27}{\%} من قيمة البر الرئيس؛ والقطيع: \qty{370}{L} في الجيل من
البقرات وحدها، و\qty{570}{L} مع ثيران منتقاة.
\end{solution}
